\chapter*{Abstract}
\addcontentsline{toc}{chapter}{Abstract}

As artificial intelligence moves toward decentralized deployment across
local edge and Dew computing nodes, preserving representational quality
without relying on cloud-based computation remains a critical challenge.
Centralized architectures impose communication latency and data-privacy
risks, whereas conventional lightweight models are prone to
representation degradation---and, in the non-generative case, to latent
collapse---when trained on resource-constrained micro-architectures.

This thesis introduces \textbf{EdgeVerify}, a Joint-Embedding Predictive
Architecture (JEPA) designed for localized execution across the Dew and
Edge computing tiers. Rather than reconstructing high-dimensional token
sequences or raw pixels, EdgeVerify predicts context representations
directly within an abstract joint-embedding space, using an
Exponential-Moving-Average (EMA) target encoder and an action-conditioned
latent predictor. To prevent dimensional collapse during localized,
layer-wise gradient updates, we integrate VICReg-inspired variance and
covariance regularization directly into the distillation objective,
removing the need for negative-sample pairing.

We present a proof-of-concept implementation of the architecture and loss
formulation, together with a reproducible experimental protocol for
measuring convergence stability, representation collapse (embedding
standard deviation and effective rank), and inference cost. A
proof-of-concept evaluation against a non-regularized baseline---conducted
on structured, digits, and CIFAR-10 data, and reproduced across two
independent environments---shows that the variance constraint reliably
prevents variance collapse (sustaining a mean embedding standard deviation
above the hinge threshold), while revealing that at the current covariance
weight the objective does not prevent dimensional (rank) collapse; this
pattern holds on the CIFAR-10 natural-image subset as well as the smaller
datasets. A systematic placement study isolates the cause: applying the
variance penalty to an embedding \emph{without} a co-located covariance
penalty on the same embedding induces rank collapse that is worse than no
regularization and invisible to both the training loss and the variance
criterion. The implied correction---co-locating both penalties on the context
embedding---restores effective rank by more than an order of magnitude on all
three datasets and recovers linear-probe accuracy to the unregularized level
(digits $0.56\!\rightarrow\!0.94$; CIFAR-10 $0.31\!\rightarrow\!0.41$).
In a federated (FedAvg) setting the correction's collapse resistance
transfers, and it holds under non-IID label skew with up to ten clients,
whereas the flawed placement collapses in every regime. Critically, however, a
larger-scale study (ResNet-18 on full CIFAR-10) shows the extreme collapse to
be largely a small-model artifact and---more consequentially---that effective
rank ceases to predict downstream accuracy at scale: the highest-rank
configuration is the \emph{least} useful, and no regularizer beats an
unregularized baseline. The rank--usefulness correspondence, and the practical
benefit of the reformulation, are therefore \emph{scale-dependent}. The
architecture itself is compact ($\sim\!0.75$M parameters, $2.84$~MB);
characterizing the scale behavior with longer, multi-seed training, along with
on-device measurement, is the primary item of future work.
